\capitulofrontal{O ENEM e a Teoria de Resposta ao Item}{tri}

A prova de Matemática e suas Tecnologias tem \textbf{45 questões} de múltipla escolha e é
aplicada no segundo dia do exame, junto com as 45 de Ciências da Natureza, em \textbf{5 horas}.
A sua nota, porém, \emph{não} é o número de acertos: ela é calculada pela
\textbf{Teoria de Resposta ao Item (TRI)}. Entender a TRI muda a forma de estudar e de fazer a prova.

\secaofrontal{O que é a TRI}

Na TRI, cada questão (item) é descrita por três números, estimados pelo INEP a partir das
respostas de milhões de participantes, e cada estudante tem uma \textbf{proficiência} $\theta$.
A probabilidade de um estudante com proficiência $\theta$ acertar o item é modelada pelo
\textbf{modelo logístico de três parâmetros}:

\begin{formula}[Modelo logístico de 3 parâmetros]
$\displaystyle P(\theta) \;=\; c \;+\; \frac{1-c}{1+e^{-a\,(\theta-b)}}$
\end{formula}

{\small\color{ebookMuted}(Em algumas apresentações aparece uma constante de escala multiplicando $a$;
ela não altera a interpretação dos parâmetros.)}

\begin{center}
\begin{tikzpicture}
\begin{axis}[width=0.9\linewidth,height=6.6cm,xmin=300,xmax=1000,ymin=0,ymax=1.02,
  xlabel={Proficiência do estudante (escala de notas do ENEM)},ylabel={Probabilidade de acerto},
  grid=major,grid style={ebookLine!60},axis lines=left,tick label style={font=\scriptsize},
  label style={font=\small},
  legend style={font=\scriptsize,at={(0.02,0.98)},anchor=north west,draw=ebookLine},
  yticklabel={\pgfmathprintnumber[fixed,precision=1,use comma]{\tick}},
  xticklabel style={/pgf/number format/1000 sep={}}]
  \addplot[nivelA,line width=1.6pt,domain=300:1000,samples=120]
    {0.20+0.80/(1+exp(-2.5*((x-500)/100-0.5)))};
  \addlegendentry{item fácil: $a=2{,}5$; $b=0{,}5$; $c=0{,}20$}
  \addplot[nivelB,line width=1.6pt,domain=300:1000,samples=120]
    {0.15+0.85/(1+exp(-2.0*((x-500)/100-1.5)))};
  \addlegendentry{item médio: $a=2{,}0$; $b=1{,}5$; $c=0{,}15$}
  \addplot[nivelD,line width=1.6pt,domain=300:1000,samples=120]
    {0.10+0.90/(1+exp(-3.0*((x-500)/100-2.8)))};
  \addlegendentry{item difícil: $a=3{,}0$; $b=2{,}8$; $c=0{,}10$}
  \draw[dashed,nivelD] (axis cs:780,0) -- (axis cs:780,0.55);
  \node[font=\scriptsize,text=nivelD,anchor=south west] at (axis cs:782,0.02) {$b=2{,}8 \Rightarrow \approx 780$};
\end{axis}
\end{tikzpicture}
\end{center}

\begin{definicao}[Os três parâmetros de cada questão]
\begin{itemize}[leftmargin=1.2em,itemsep=2pt,topsep=2pt]
  \item[$b$] \textbf{Dificuldade.} Indica em que ponto da escala a curva “sobe”: é a proficiência
        para a qual a chance de acerto fica a meio caminho entre $c$ e 100\%. Quanto maior o $b$,
        mais difícil o item.
  \item[$a$] \textbf{Discriminação.} É a inclinação da curva. Um $a$ alto significa que o item separa
        bem quem domina o assunto de quem não domina.
  \item[$c$] \textbf{Acerto ao acaso.} É a chance de acerto de quem não sabe nada (o “chute”).
        Com cinco alternativas fica perto de $0{,}2$, mas é estimado item a item.
\end{itemize}
\end{definicao}

\secaofrontal{Da TRI para a nota}

A escala de notas do ENEM foi construída de modo que a proficiência $\theta = 0$ corresponda a
\textbf{500 pontos} e cada desvio-padrão corresponda a \textbf{100 pontos}. Por isso, nesta coleção,
mostramos a posição aproximada de cada questão na escala como $500 + 100b$: uma questão com
$b = 1{,}5$ fica por volta de 650 pontos.

O gráfico abaixo mostra onde estão as \textbf{\totalbanco\ questões} de Matemática válidas
(não anuladas) aplicadas no caderno azul entre \anosbanco. A mediana fica em torno de
\textbf{\trimediana\ pontos}: metade das questões de Matemática exige um desempenho acima
disso. É por isso que acertar as questões fáceis e médias é tão valioso.

\begin{center}
\begin{tikzpicture}
\begin{axis}[ybar,bar width=8.5pt,width=0.92\linewidth,height=5.6cm,
  xmin=300,xmax=1050,ymin=0,axis lines=left,
  xlabel={Posição da questão na escala ($500+100b$)},ylabel={Nº de questões},
  tick label style={font=\scriptsize},label style={font=\small},
  xtick={300,400,...,1000},xticklabel style={/pgf/number format/1000 sep={}},
  ymajorgrids,grid style={ebookLine!60}]
  \addplot[fill=ebookPrim,draw=none] coordinates {\trihist};
  \draw[dashed,ebookAcc,line width=1pt] (axis cs:\trimediana,0) -- (axis cs:\trimediana,175)
    node[above,font=\scriptsize,text=ebookAcc] {mediana $\approx$ \trimediana};
\end{axis}
\end{tikzpicture}
\end{center}

Pelos critérios desta coleção, essas questões se distribuem assim:
\textcolor{nivelA}{\faSquare}\,\textbf{\trinA} fáceis,
\textcolor{nivelB}{\faSquare}\,\textbf{\trinB} médias,
\textcolor{nivelC}{\faSquare}\,\textbf{\trinC} difíceis e
\textcolor{nivelD}{\faSquare}\,\textbf{\trinD} muito difíceis.

\secaofrontal{Coerência pedagógica: por que o mesmo número de acertos dá notas diferentes}

A TRI procura a proficiência que melhor explica o \emph{conjunto} das suas respostas. Quem acerta
questões difíceis mas erra muitas fáceis apresenta um padrão pouco provável para alguém que domina o
conteúdo — um padrão típico de chute. Veja um exemplo \emph{ilustrativo}:

\begin{center}
\renewcommand\arraystretch{1.35}
\begin{tabular}{@{}lcccc@{}}
\toprule
& \textbf{Fáceis} (de 15) & \textbf{Médias} (de 15) & \textbf{Difíceis} (de 15) & \textbf{Total} \\
\midrule
Ana   & 14 & 9 & 2 & 25 acertos \\
Bruno & 6  & 9 & 10 & 25 acertos \\
\bottomrule
\end{tabular}
\end{center}

Ana e Bruno acertaram 25 questões, mas a nota de Ana tende a ser \textbf{maior}: o padrão dela é
coerente (acerta o fácil, erra parte do difícil). Os acertos de Bruno nas difíceis “valem menos”,
porque a TRI considera a chance de que tenham vindo do acaso, e os erros nas fáceis pesam contra ele.

\secaofrontal{Estratégias de prova com a TRI}

\begin{dica}[Garanta as fáceis]
Errar uma questão fácil custa caro. Leia com calma as questões curtas e diretas e revise as contas
delas antes de investir tempo nas mais longas.
\end{dica}
\begin{dica}[Faça a prova em duas voltas]
Na primeira volta, resolva tudo o que for rápido e marque as questões trabalhosas. Na segunda, volte
a elas com o tempo que sobrou. São 5 horas para 90 questões no segundo dia — cerca de
\textbf{3 minutos por questão}, já descontado o tempo de preencher o cartão-resposta.
\end{dica}
\begin{macete}[Nunca deixe em branco]
Questão em branco conta como erro. Se não souber, \textbf{elimine alternativas} absurdas e marque
uma das restantes: eliminar duas alternativas eleva a chance de acerto de 20\% para cerca de 33\%.
\end{macete}
\begin{atencao}[Chute não substitui estudo]
Acertos ao acaso em questões difíceis ajudam pouco quando há erros nas fáceis. A melhor estratégia
de TRI é dominar o básico de \emph{todos} os conteúdos — exatamente a proposta desta coleção.
\end{atencao}

\secaofrontal{Como esta coleção usa a TRI}

\begin{itemize}[leftmargin=1.4em,itemsep=2pt]
  \item As questões oficiais vêm dos cadernos azuis do segundo dia (\anosbanco) e trazem os
        parâmetros $a$, $b$ e $c$ publicados pelo INEP nos \emph{Microdados do ENEM}
        (arquivo de itens da prova).
  \item Em cada capítulo, elas aparecem em \textbf{ordem crescente de $b$}. As inéditas seguem a
        mesma progressão, com o nível estimado a partir das oficiais do mesmo assunto.
  \item A posição $500 + 100b$ é uma aproximação útil para comparar questões; a dificuldade que
        você sente depende do seu preparo em cada assunto.
\end{itemize}
